% !PRINTED 646
\weberpage{646}
%<<CONT
% [sic] Print changes the argument list here to $\xi_0,\ldots,\xi_{n-2}$;
% the function introduced on printed p. 645 has arguments through $\xi_{n-1}$.
2. but neither does this function change when the cyclic permutation $(k_0, k_1 \ldots k_{n-2})$ is performed among the $k_\nu$, and since now the quantities $k_\nu$ are the roots of a cyclic equation in the field $\mathfrak{K}$, the function $S(\xi_0, \xi_1 \ldots \xi_{n-2})$ is a quantity in $\mathfrak{K}$, i.e. rational. If one applies this to the coefficients of the equation whose roots are the quantities (10), it follows that these quantities are the roots of an equation of degree $n$ in $\mathfrak{K}$.

As for the irreducibility of this equation, the following is to be remarked on that point:

The radicals $\sqrt[n]{R_\nu}$ are, if one of them is rational, all rational, and if one is irrational, so are the others [§.\,185, (8)]. Which of the two cases occurs therefore depends, by §.\,179, on whether the function $x^n - R_0$ is reducible or irreducible in $\mathfrak{K}$. If it is reducible, then one of the quantities $\xi_h$ is rational in the field $\mathfrak{K}$, and the equation for $\xi$ is reducible. All the quantities $\xi$ are contained in the field $\mathfrak{K}(\varepsilon)$, whose degree is at most equal to $n-1$.

But if $\sqrt[n]{R_0}$ is irrational, then $x^n - R_0$ is irreducible in $\mathfrak{K}$. Every equation $\Phi(\xi_0) = 0$ passes, on substitution of the value of $\xi_0$, into an equation $\Psi\!\left(\sqrt[n]{R_0}\right) = 0$ whose coefficients are contained in $\mathfrak{K}$, and which must therefore continue to hold when $\sqrt[n]{R_0}$ is replaced by any one of its $n$ values.

But thereby $\xi_0$ passes into each of the values $\xi_0, \xi_1 \ldots \xi_{n-1}$, i.e. the equation of degree $n$ whose roots are the $n$ quantities $\xi$ is irreducible in $\mathfrak{K}$.

We can therefore state the converse of Theorem I as follows:

\gesperrt{II. Every quantity $\xi$ contained in the form (9) is the root of an equation of degree $n$ in $\mathfrak{K}$, and this equation is irreducible if one leaves aside the special case that one of its roots is rational and the others are contained in $\mathfrak{K}(\varepsilon)$.}
